% Standalone fragment.  Requires amsmath, amssymb, amsthm and hyperref.
% Proof relative to the explicit analytic lemmas of the original manuscript.
\subsection{A weighted moment for the reconstructed physical numerator}
\label{sec:new-weighted-numerator}

The following is a new analytic amendment to the high estimate.  Its
inputs are the exact local identities, smooth reflection and scale
supremum lemmas, and the fourth moment with short prime factors
\textup{(Lemma \texttt{lem:plain})} of the original manuscript.  Those
inputs are imported, rather than reproved here. The coefficient bounds,
literal Mellin reconstruction, original-numerator fourth moment, and
their composition over the amplitude bins are checked in Lean
(Section~\ref{sec:lean-correspondence}); the composition is
\nolinkurl{WeightedQRH.Numerator.source_integrated_amplitude_classes},
on which \nolinkurl{WeightedQRH.Numerator.actual_raw_fullW_input} is
built. The argument below requires coefficient estimates for the Euler
correction, not merely pointwise bounds for an Euler product.

Write $z_0=17/50$, $\kappa=3/4$, let $\theta$ denote the target exponent
($\theta=B$ in Section~\ref{sec:new-certificate}), and assume the
already proved global bound $\beta_*\le7/8$.  Thus the zero-free premise of the imported plain
moment holds.  This choice of $\kappa$ is fixed throughout and uses no
feedback improvement.

\paragraph{The estimate.}
Let $U=Z^d$, $Y=Z^{l_y}$, and let $\mathcal B$ be a fixed dynamic bin of
nonprincipal physical rows $u$ with $q_u\asymp U$.  Put
$\delta=2a-1$, where $51/100\le a\le7/8$. Let $\ell$ be the total physical
prime-slot length.  Suppose
\begin{equation}\label{eq:new-capacity-supply}
 \frac d2<l_y<d,\qquad
 0<\frac{2l_y-d}{6\kappa}\le\ell.
\end{equation}
For each pointwise amplitude and witness subdivision $\mathcal C$ of
this bin, suppose the original row-count argument gives
\[
 \#\mathcal C\ll U^{R+\epsilon}(1+T_1)^A,
 \qquad 0\le q\le\delta/2,
\]
where $q\ell$ is the sum of slot length times amplitude, assigning
amplitude zero to every error slot.  Up to arbitrarily small exponent
losses and a fixed polynomial in the retained heights, its contribution
to the high comparison is bounded by the supremum of
\begin{equation}\label{eq:new-weighted-high-exponent}
 \begin{split}
 E_W(d)={}&a-\theta+h(z_0-1/6)-\frac{l_y}{2}-\frac\ell2+q\ell\\
 &+d\left(\frac{3R}{4}+\frac14-z_0\right)
       -\frac{q(2l_y-d)}{12\kappa}.
 \end{split}
\end{equation}
The supremum is taken after continuing the whole fixed dynamic bin.
Individual amplitude sets are not continued separately.  In particular,
choosing between this estimate and the original high estimate for a
whole bin gives
$\min\{\sup_q E_{\rm old},\sup_q E_W\}$, without requiring an
interchange of the minimum and supremum.

\paragraph{Coefficient norms and the local correction.}
Fix the retained $s,z$ with $\Re s=a+16e$ and $\Re z=z_0$.
Set
\[
 \rho=\delta/4,\qquad \sigma_-=1/2-\rho.
\]
Then $\sigma_->0$ and
$\Re s+\sigma_--1\ge\delta/4>0$.  Both $\Re w=1/2$ and
$\Re w=\sigma_-$ lie in the first Euler region.  For a polynomial in
$Q^{-w}$ define
\[
 \left\|\sum_j b_jQ^{-jw}\right\|_\sigma
       =\sum_j|b_j|Q^{-j\sigma}.
\]
The local denominators $1-D$, $1-R$, and $1-V$ in the original
manuscript are independent of $w$ and uniformly separated from zero.
The factor $H_p$ is a polynomial of degree at most two in $Q^{-w}$ at
$p\nmid u$, and at most one at $p\mid u$.  The exact defect identity
\[
 H_p-1=
 \frac{D(V+W-VW)-VW+(1-V)(1-W)\mathcal E_p}{1-D},
 \qquad \mathcal E_p=P_p^*+D,
\]
therefore gives, by taking absolute coefficients \emph{after} this
cancellation,
\[
 \|H_p-1\|_{\sigma_-}\ll
 \begin{cases}
 Q^{-1-\epsilon_H},&p\nmid u,\\
 Q^{-\epsilon_H},&p\mid u,
 \end{cases}
\]
for a fixed $\epsilon_H>0$.  The same holds at $\sigma=1/2$.
The convergent prime product and the divisor-product estimate imply
that every unselected product, excluding any tuple of distinct slot
labels, has absolute coefficient norm $\ll U^\epsilon$ on both lines.
These assertions are uniform in the retained imaginary parts.

Write
$E_p=G_p+\overline{\chi_p}(u)H_p$.  For an error-slot subset $I$,
the quotient-free part of the correction is
\begin{equation}\label{eq:new-error-coefficient-product}
 F_I(w)=\sum_{(p_i)_{i\in I}}
 \prod_{i\in I}\left[
 W_i(q_{p_i}/P_i)q_{p_i}^{z-1}E_{p_i}\right]
 \prod_{\substack{p\notin S\\p\notin\{p_i:i\in I\}}}H_p.
\end{equation}
All remaining slots retain their actual main factors
$-P_i^{z-1/2}Q_i(u;z)$.  Equation
\eqref{eq:new-error-coefficient-product} follows by expanding the
holomorphic selected-tuple identity.  It does not divide by $H_p$ and
does not alter a main-slot mask or coefficient.

At $p\nmid u$, the exact local error cancellation
\textup{(original equation \texttt{old-eq:5.13b})} yields
\[
 \|E_p\|_\sigma\ll
 Q^{-s_r}+Q^{-6z_0}+Q^{4-5s_r-6z_0}+Q^{1-\sigma-6z_0}
 \qquad(\sigma\in\{1/2,\sigma_-\}).
\]
Each exponent is at most $-1/2$ with a fixed margin.  The ideal count
therefore bounds the sum over a slot, including $Q^{z-1}$, by
$O(P_i^{z_0-1/2})$.

At $p\mid u$, with $j=v_p(u)$, the main factor is zero and
\[
 E_p=G_p=\overline{\eta(p)}Q^s(1-V)\mathcal E_p-Q^{-w}H_p.
\]
After dividing the slot factor by $Q^{z_0-1/2}$, the non-tail boundary
terms of the six-valuation table have the following exponents on
$\Re w=\sigma_-$, at $e=0$:
\[
\begin{array}[b]{c|ccccc}
 j&1&2&3&4&5\\ \hline
 \text{exponents}&-1+\rho&-\delta&-\delta+\rho\ \text{and}\ {-3\delta/2}
   &-3\delta/2+\rho&-5\delta/2
\end{array}.
\]
All are negative; positive $e$ improves them.  The $R,V$ tails and
$-Q^{-w}H_p$ are smaller.  The only exception is the strict term
\[
 -\eta(p)(Q-1)Q^{-s-w},\qquad j\ge2,
\]
whose normalized slot factor costs at most $Q^\rho$ on the left line
and at most one on the central line.  Ramified labels are divisor-many.
For any fixed triangle-expanded tuple, let $J$ consist precisely of
its strict ramified labels.  They are distinct.  Choose the original
small-presentation character, whose ideal coefficients agree with the
physical row. Let $C_u$ be the product of its primitive conductor norm
and full redundant Euler-ideal norm. This combined scale satisfies
the original conductor allocation
\[
 C_u\le C_S\operatorname{rad}(u)
 \le C_*U\prod_{p\in J}Q^{-(j-1)}.
\]
Consequently
$(C_u/(C_*U))^\rho\prod_{p\in J}Q^\rho\le1$.
The single conductor deficit is applied to the whole tuple.

Expand $F_I(w)=\sum_k A_{I,u}(k)q_k^{-w}$ and put
$B_I=\prod_{i\in I}P_i^{z_0-1/2}$.  The preceding estimates prove
\begin{align}
 \sum_k|A_{I,u}(k)|q_k^{-1/2}&\ll U^\epsilon B_I,
 \label{eq:new-central-coefficient-norm}\\
 \left(\frac{C_u}{C_*U}\right)^\rho
 \sum_k|A_{I,u}(k)|q_k^{-1/2+\rho}&\ll U^\epsilon B_I.
 \label{eq:new-left-coefficient-norm}
\end{align}
Combining them gives the useful weighted form
\begin{equation}\label{eq:new-weighted-coefficient-norm}
 \sum_k|A_{I,u}(k)|q_k^{-1/2}
 \max\left(1,\frac{C_uq_k}{C_*U}\right)^\rho
 \ll U^\epsilon B_I.
\end{equation}
This also explains why a strict ramified error does not consume prime
capacity: its forced dilation is paid for by a distinct factor in the
conductor deficit.

\paragraph{Reconstructing and reflecting the numerator.}
For each fixed nonprincipal physical row, move the complete $w$
integral to $\Re w=1/2$.  The numerator is entire; the correction is
holomorphic throughout this move, and the Mellin profile has arbitrary
decay.  The bin itself is independent of $w$.  The absolute coefficient
norm just proved and polynomial strip growth justify termwise Mellin
inversion:
\begin{equation}\label{eq:new-exact-w-reconstruction}
 \begin{split}
 &\frac1{2\pi i}\int_{(1/2)}
 Y^{w-1}\widehat W_1(w)L^S(w,\chi_u)F_I(w)\,dw\\
 &\hspace{8mm}=
 Y^{-1/2}\sum_k A_{I,u}(k)q_k^{-1/2}
              S_{\chi_u}(Y/q_k;W_1).
 \end{split}
\end{equation}
Here $S$ denotes the original centrally normalized plain sum.  Since
$W_1$ has compact support, say in $[r_-,r_+]$, terms with $q_k>r_+Y$
vanish identically before reflection.  Thus all relevant scales lie
in a fixed polynomial range.

The imported smooth functional equation replaces $S_{\chi_u}(X;W_1)$ by reflected sums
with conjugate character and fixed transformed profiles. Restoring
natural zeros follows the original reflection proof. The redundant
Euler restoration need not be supported on fixed $S$: every restoration
subset has norm at most the full redundant ideal norm, and every
smooth-ideal denominator decreases the reflected scale. Consequently
all reflected scales are at most $C_u/X$, with the combined $C_u$
defined above. The original restoration coefficient mass is subpower
in the small modulus. Fixed constants enter $C_*$. No extra weighted
restoration mass or new moving mask is needed.

Use \eqref{eq:new-weighted-coefficient-norm} to remove all row-dependent
coefficients before applying a moment.  It leaves a rowwise supremum
of reflected sums of scale at most
$C_*(U/Y)T$, multiplied by $T^{-\rho}$, where
\[
 T=\max(1,C_uq_k/(C_*U)).
\]
The imported smooth scale-supremum lemma handles the row-dependent
scale using logarithmically many boxes and finitely many derivatives.
Annular decomposition of the reflected profile gives the usual
arbitrarily small length loss, summable lower annuli and arbitrarily
decaying upper tails.  Since $q_k\le r_+Y$, all main reflected scales
are $O(U)$ and the declared lengths stay bounded.  Thus no uniform
moment assertion for unbounded lengths is being used.

\begin{proposition}[Checked original-numerator moment]\label{prop:checked-original-fourth}
Assume $51/100\le\beta_*$. \textup{(}The bound $\beta_*\le7/8$ of
\cite{openai} is available throughout but is not needed here; the Lean
statement assumes only the lower bound.\textup{)}
Fix the original weighted source data, its moment loss $t>0$, a
Schwartz profile $G$, and positive constants $\lambda,\xi,b$.
For sufficiently large $U$, let $\mathcal R$ be any finite set of
nonprincipal physical rows with $U^{1/100}\le q_u\le U$.
For every row choose $0<N_u\le bY$, where $Y>0$, and its actual
combined conductor $C_u$. Suppose, for a common $m\ge0$,
\[
 m+\xi\le2,\qquad
 C_u\le C_S U^m\frac{Y}{N_u}.
\]
Let $T$ be any selected family of original prime slots with nonnegative
widths $r_i$ below the certified mesh
$\operatorname{fineMesh}(2,0,1,3/4,t/4)$, and let
$\Re z_i=17/50$, $|\Im z_i|\le H$, $H\ge0$. Assume
\[
 2(m+\xi)+\frac92\sum_{i\in T}r_i\le1
 \quad\hbox{or}\quad T=\varnothing,
 \qquad
 \left|\prod_{i\in T}Q_i(u;z_i)\right|\le U.
\]
Then there are a fixed integer $J\ge0$ and $C>0$ such that
\[
 \sum_{u\in\mathcal R}
 \left|S_{\chi_u}(Y/N_u;G)\right|^4
 \left|\prod_{i\in T}Q_i(u;z_i)\right|^2
 \le C(1+H)^J U^{1+t+\lambda}.
\]
The constants and threshold precede all choices of rows, scales,
selected slots, and heights. They may depend on the fixed source data,
profile and displayed fixed positive constants.
\end{proposition}
\begin{proof}
Use the actual smooth reflection identity with unit-modulus phase and
cutoff $U^{\xi/4}$. The certified retained-profile theorem bounds the
weighted fourth moment of the retained part by
$C_r(1+H)^J U^{1+t+\lambda}$. It uses the positive-capacity source
moment when $T\ne\varnothing$ and the unrestricted bounded-length
zero-slot source moment otherwise. Both applications preserve the
original physical coefficients and selected factors.

For the discarded part, $C_u\le C_S q_u\le C_SU$ and $N_u\le bY$
give a uniform reflected-scale bound $C_uN_u/Y\le bC_SU$.
The original arbitrary-order discarded-tail estimate therefore gives
any power saving, uniformly over these rowwise choices. The elementary
row count is at most $128U$, and the assumed selected product is at
most $U$. Choosing sufficient tail decay bounds its weighted fourth
moment by a constant $C_d$. Finally,
$|r+d|^4\le8(|r|^4+|d|^4)$, the reflection phase has modulus one,
and $(1+H)^J U^{1+t+\lambda}\ge1$. The result follows with
$C=8(C_r+C_d)$. The Lean theorem allows the more general fixed bound
$C_{\rm slot}U^{q_{\rm tail}}$ on the selected product; the displayed
statement specializes both constants to one.
\end{proof}

\paragraph{Selecting prime factors inside the fourth moment.}
Fix an amplitude set $\mathcal C$ after the full $w$ integration, at
one retained $(s,z)$ and witness tuple.  Put
\[
 L=\ell/d,\qquad n_0=1-l_y/d,
 \qquad z_P=\frac{1-2n_0}{6\kappa}.
\]
For a scale dyad $T=U^v$, allow length $n(v)=n_0+v$ and define
\[
 z(v)=\max\left(0,z_P-\frac{2v}{6\kappa}\right).
\]
The supply condition gives $z_P\le L$.  Select the largest positive
amplitude main slots up to this capacity.  If their total length is
smaller, select all of them.  The original selection lemma, with a
sufficiently fine mesh, gives
\[
 |Q_{\rm selected}|\ge U^{qz(v)-o(1)},
 \qquad |Q_{\rm all}|\le U^{qL+o(1)}
\]
on $\mathcal C$.  Actual selected length never exceeds $z(v)$; no
zero-amplitude factor is required for this lower bound.  The reflected
numerator has character $\overline{\chi_\bullet}(u)$, exactly matching
the physical main slots $\overline{\chi_p}(u)$.  Conjugating the whole
plain/selected-slot product inside its absolute value gives the positive
row orientation in the statement of the imported plain moment; it only
conjugates the profiles and fixed finite-ray coefficients, and preserves
every zero extension.  For sufficiently large $U$ in the moderate
dyads, a physical row cannot induce a character in the fixed exceptional
group $\Theta$, since a good prime of that row ramifies.  The finitely
many unit or fixed-prime-supported exceptions belong to the separate
small-row estimate.

H\"older's inequality and the imported plain moment give
\begin{align*}
 \sum_{u\in\mathcal C}|S_{\bar\chi_u}(n(v))Q_{\rm all}(u)|
 &\ll U^{qL-qz(v)/2+o(1)}(\#\mathcal C)^{3/4}
 \left(\sum_u |S_{\bar\chi_u}(n(v))|^4
                     |Q_{\rm selected}(u)|^2\right)^{1/4}\\
 &\ll U^{qL-qz(v)/2+3R/4+1/4+o(1)}.
\end{align*}
If $z(v)>0$, the capacity equation is $2n(v)+6\kappa z(v)=1$.
If $z(v)=0$, use the original lemma's stronger zero-slot assertion:
it gives $U^{1+o(1)}$ for bounded nonnegative plain lengths, without
a length restriction. This is the certified \texttt{ZeroAt} input,
not the weaker downstream detector estimate with exponent
$\max(1,2n)$. Small slack is reserved for annuli and scale boxes.
The derivative versions used by the scale supremum have the same
exponent and a fixed height cost.

The coefficient weight adds $-\rho v$.  Before the capacity is
exhausted, the exponent relative to $v=0$ is
\[
 \left(\frac q{6\kappa}-\rho\right)v
 \le-\frac{5\delta}{36}v,
\]
because $\kappa=3/4$, $q\le\delta/2$, and $\rho=\delta/4$.
After capacity is exhausted, its slope is $-\rho$. Thus every
dilation dyad is controlled by $v=0$, with room for the small fixed
losses. All error subsets satisfy the same estimate; no lower bound
for an error factor is used.

Restoring the main-slot normalizers, the $w$-integrated bin is bounded
by
\[
 Y^{-1/2}Z^{\ell(z_0-1/2)+q\ell}
 U^{3R/4+1/4-qz_P/2+o(1)}.
\]
The remaining $q_u^{-z}$ factor, reciprocal-denominator estimate and
$X,Z$ Mellin weights are those of the original high-bin proof.
They give exactly \eqref{eq:new-weighted-high-exponent}.  The finite
amplitude, witness and error subdivisions are used only in this
estimate of the already continued whole bin.  Rapid external Mellin
decay and the fixed height orders treat the joins and discarded
coordinates with the original cumulative height allocation.

\paragraph{Dependence on the row scale.}
In the positive-capacity range the coefficient of $d$ in
\eqref{eq:new-weighted-high-exponent} is
\[
 \frac{3R}{4}+\frac14-z_0+\frac q{12\kappa}.
\]
For the original selected count $R\ge1-\delta$ and
$\delta\le5/6$, this is at least $7/200>0$.  Thus $d=h$ controls
the moderate row range whenever \eqref{eq:new-capacity-supply} holds
throughout that range.  The small and far row estimates can remain
those of the original proof.

Finally, at fixed $R,q$ the difference from the original high exponent
is
\[
 E_W-E_{\rm old}
 =\frac{\delta l_y}{2}
   +d\left(\frac{1-R}{4}-\frac\delta2\right)
   -\frac{q(2l_y-d)}{12\kappa}.
\]
The final term is the additional saving supplied by the reconstructed
physical numerator and the existing prime-weighted fourth moment.
It is independent of any parameter-feedback argument.
